\section{From completed perfections to ordinary cochains}
\label{sec:decompletion}

The geometric calculation concerns the completed perfections \(B^\flat\)
of finite marking stages \(B=L^U\), whereas ordinary cooperations contain
\(L\). Theorem~\ref{thm:decompletion} first compares the root-stage union
\(R=L^{\mathrm{perf}}\) with the completed-stage cochains.
Theorem~\ref{thm:coefficient-calculation} then uses a distribution operator
and the order-two norm quotient to compare \(L\) with \(R\) on
\(P\)-cochains and compute the four ordinary cohomology lines. The central
input is finite-stage finiteness, proved in
Proposition~\ref{prop:finite-stage-finiteness}.

Throughout this section cohomology of \(L\) and \(R\) uses
Definition~\ref{def:stage-cochains}: a cochain with a profinite parameter
takes its values in one complete finite marking and root stage. Write
\(C_{\mathrm{cts}}(T,M)\) for continuous functions with the uniform topology
when \(T\) is compact. For complete coefficients,
\[
 C_{\mathrm{cts}}\bigl(T,C^r_{\mathrm{cts}}(G,M)\bigr)
 =C_{\mathrm{cts}}(T\times G^r,M).
\]
For stage unions, the corresponding expression means the colimit of these
fixed-stage spaces.

\subsection{Compact duality and the coefficient topology}

The finite resolution of Section~\ref{sec:finite-resolutions} now supplies
compact duality. We then establish the strictness and parameter-lifting
facts needed to return from completed coefficients to ordinary cochains.

\begin{lemma}[Compact--discrete duality]
\label{lem:compact-duality}
Let $M$ be a compact continuous $\Fp$-representation of $H$ and give
\[
 M^\vee=\Hom_{\Fp,\mathrm{cts}}(M,\Fp)
\]
the discrete topology and the contragredient action.  Continuous
$H$-cohomology of $M$ is compact Hausdorff in its ordinary cochain
quotient topology, and there are natural isomorphisms
\begin{equation}
 \label{eq:compact-cohomology-duality}
 H^a_{\mathrm{cts}}(H,M)^\vee
 \cong H^{8-a}_{\mathrm{cts}}(H,M^\vee).
\end{equation}
The transposed coefficient map induces the dual cohomology map.
The reversed assertion holds for a discrete module and its compact
dual.
\end{lemma}

\begin{proof}
Use the finite resolution in Lemma~\ref{lem:finite-resolution}.
Its Hom complex on $M$ has compact terms and closed boundaries, so its
cohomology is compact Hausdorff.  The continuous comparisons with bar
cochains identify the resulting topology with the bar-cochain quotient
topology.  This does not assert compactness of the bar-cochain spaces.

Continuous $\Fp$-duality is exact on compact vector spaces: a
functional on a closed subspace extends after passing to a suitable
finite-dimensional quotient.  If $Q$ is finite projective over
$\Lambda=\Fp[[H]]$, put $Q^*=\Hom_\Lambda(Q,\Lambda)$ and turn this
right module into a left module using the involution $h\mapsto h^{-1}$.
Finite-projective evaluation gives
\[
 \Hom_{\Lambda,\mathrm{cts}}(Q,M)^\vee
 \cong \Hom_{\Lambda,\mathrm{cts}}(Q^*,M^\vee).
\]
For a finite free module this is coordinatewise evaluation; taking
projective summands proves the general case.  The identification is
compatible with precomposition and coefficient maps.  Apply it to
each term, reverse degrees about eight, and use the reversed dual
resolution in Lemma~\ref{lem:finite-resolution}.  Exactness of compact
duality proves \eqref{eq:compact-cohomology-duality} and its naturality.
Dualizing again proves the discrete assertion.
\end{proof}

\begin{lemma}[Strictness]
\label{lem:strictness}
Let $A^\bullet$ be a complex of complete, separable, metrizable,
linearly topologized $\Fp$-spaces with continuous linear differentials.
If $H^a(A^\bullet)$ is finite, the boundary subgroup is
open and closed in the cycle group, and the differential onto that
boundary subgroup is open.  Thus, for every neighborhood $V$ of zero
in $A^{a-1}$, sufficiently small $a$-cocycles have primitives in $V$.
\end{lemma}

\begin{proof}
Put $E=A^{a-1}$, $Z=\ker d^a$, and $f=d^{a-1}:E\to Z$.
The closed subspace $Z$ is complete and metrizable.  Choose decreasing
bases of open linear subspaces $E=E_0\supset E_1\supset\cdots$ in
$E$ and $W_0\supset W_1\supset\cdots$ in $Z$.
Separability makes $E/E_n$ countable.  Since $Z/f(E)$ is finite,
$Z/f(E_n)$ is countable as well.  Set
\[
 H_n=\overline{f(E_n)}\subset Z.
\]
Countably many cosets of this closed subgroup cover $Z$.
Baire's theorem gives $H_n$ nonempty interior, so it is open.

We show that $f(E_n)=H_n$ for every $n$.
Given $z\in H_n$, put $r_n=z$.  If $r_j\in H_j$, density of
$f(E_j)$ in $H_j$ and openness of $H_{j+1}\cap W_{j+1}$ allow a
choice $e_j\in E_j$ with
\[
 r_{j+1}=r_j-f(e_j)\in H_{j+1}\cap W_{j+1}.
\]
Indeed the translated neighborhood of $r_j$ lies in $H_j$ and meets
the dense image.  The series $\sum_{j\ge n}e_j$ is Cauchy because
its tails lie in the decreasing open subspaces $E_j$.
Completeness gives its sum $e\in E_n$, since $E_n$ is closed.
The residuals tend to zero, so continuity gives $f(e)=z$.
Thus every $f(E_n)=H_n$ is open.  Taking $n=0$ proves that the
boundary subgroup is open and closed in $Z$; the other values of
$n$ prove that the differential onto it is open.  Choosing $E_n$
inside a prescribed neighborhood gives the last assertion.
\end{proof}

\begin{lemma}[Continuous profinite parameters]
\label{lem:profinite-parameters}
Let $T$ be an arbitrary compact profinite space.
\begin{enumerate}
\item A continuous map from $T$ to $F$ lifts across an open continuous
surjection $E\twoheadrightarrow F$ of complete metrizable linearly
topologized groups.
\item If a complex $A^\bullet$ of complete, separable, metrizable, linearly
topologized $\Fp$-spaces with continuous linear differentials has finite
cohomology, then naturally
\begin{equation}
 \label{eq:finite-cohomology-parameters}
 H^a C_{\mathrm{cts}}(T,A^\bullet)
 \cong C_{\mathrm{cts}}(T,H^a(A^\bullet)),
\end{equation}
where the group on the right is discrete.
\item For a filtered diagram of such complexes, this identity commutes
with the cochain colimit, with the colimit of cohomology groups given
the discrete topology.
\end{enumerate}
\end{lemma}

\begin{proof}
Choose a decreasing basis of open subgroups $E_n$ of $E$ tending to
zero, small enough that their images also tend to zero in $F$.
Openness makes each $q(E_n)$ open.  Approximate the given function on
a finite clopen partition of $T$ by finitely many values and lift those
values to $E$, leaving an error in $q(E_1)$.  Refine the partition and
approximate that error by values with lifts in $E_1$, leaving an error
in $q(E_2)$.  Continue.  The corrections form a uniformly Cauchy
sequence because all later corrections lie in $E_n$; completeness
produces a continuous lift.  No metrizability of $T$ is used.

For the second assertion, a continuous family of cycles has a
continuous image in the finite, discrete cohomology group by
Lemma~\ref{lem:strictness}.  A family mapping to zero lifts through
the open differential onto boundaries by the first assertion.
Conversely, choose a cycle representative for each of the finitely
many values of a family of cohomology classes.  This proves
\eqref{eq:finite-cohomology-parameters}.  Finally, a continuous map
from $T$ to a discrete filtered colimit has finite image.  Its finitely
many values, and every equality needed between them, occur at one
stage.  Hence
\[
 \colim_i C_{\mathrm{cts}}(T,V_i)
 \cong C_{\mathrm{cts}}(T,\colim_i V_i)
\]
for discrete vector spaces $V_i$, even when their transition maps are
not injective.  Exactness of filtered colimits proves the last claim.
\end{proof}

We will repeatedly use two elementary consequences of the coefficient
convention.  A short exact coefficient sequence induces a short exact
sequence of continuous cochain complexes whenever its surjection has
a continuous section as a map of spaces; additivity and equivariance
of the section are unnecessary.  For a stage union, it suffices that
every target stage lift continuously into one larger source stage and
that the kernel have the induced topology.  Quotients by open
valuation lattices are discrete and admit such sections.

Also, if $G$ is compact and second countable and $D$ is countable
discrete, every $C^r_{\mathrm{cts}}(G,D)$ is countable.  Indeed $G^r$
has only countably many clopen sets, and a cochain is described by a
finite clopen partition and finitely many values.  This observation
will be used only with $T$ a point.

The finite marking stages are products of local fields with finite
residue fields.  They, their finite root stages, and their completed
perfections are Polish.  For the latter assertion, choose a
countable dense subset at each root stage; their union is dense
in the completed perfection.  Cochain spaces on the compact
metrizable spaces $H^r$ are Polish for the uniform topology:
completeness is uniform completeness, and locally constant functions
on countably many finite clopen partitions supply separability.
The valuation balls are invariant.  The same assertions hold for
closed order lattices.  Thus all applications below of
Lemmas~\ref{lem:finite-resolution} and~\ref{lem:strictness} have the
stated hypotheses.  For arbitrary $T$, we instead use
Lemma~\ref{lem:profinite-parameters}; we do not assert that its
parameterized cochain spaces are Polish or countable.

\subsection{Distributions and the invariant differential}

The normalized generator-lift torsors of Lemma~\ref{lem:lift-torsors}
are torsors under $\Gamma_2[p^l]$, with coordinate algebras
\[
 R_l=L^{1/p^{2l}},\qquad R=\colim_l R_l.
\]
Their translation action commutes with $H$.  Changes of connected
marking act on translations by continuous constant automorphisms.

\begin{lemma}[The distribution operator]
\label{lem:distribution-operator}
After the fixed finite enlargement of $k$, the inverse-limit
distribution algebra of the translation tower is $k[[D]]$.  A
parameter $D$ can be chosen with the following properties.
\begin{enumerate}
\item Its action on $R$ is $L$-linear and locally nilpotent, commutes
with $H$, and satisfies
\begin{equation}
 \label{eq:distribution-kernels}
 \ker(D^{p^i}:R\to R)=L^{1/p^i}\quad(i\ge0),
 \qquad D:R\twoheadrightarrow R.
\end{equation}
\item The order-two norm quotient acts on $D$ by its nontrivial
character.  The action of $J$ on $k[[D]]$ is continuous on finite jets.
\item If $i$ is positive, even, and divisible by $[k:\Fp]$, and
$q=p^i$, then
\begin{equation}
 \label{eq:distribution-frobenius}
 D(f^q)=(D^qf)^q,\qquad
 D^q(af)=aD^qf\quad(a\in L^{1/q}).
\end{equation}
\item Each operator on a fixed root stage is continuous after passage
to a finite larger marking stage.  Every finite target stage has a
continuous additive lift under $D$ at some finite larger marking and
root stage.  Thus the sequence with kernel $L$ and operator $D$ is
exact on the prescribed cochains, also with profinite parameters.
\end{enumerate}
\end{lemma}

\begin{proof}
The Cartier dual of the height-two, dimension-one Honda group is
again connected of height two and dimension one; see
\cite[Section~2.2, Proposition~1, and Section~2.3, Proposition~3]{Tate67}
and the Frobenius--Verschiebung calculation below. A coordinate on
that dual identifies the finite distribution algebra with
$k[D]/(D^{p^{2l}})$; the transition maps are restriction to the dual
torsion kernels.  Their inverse limit is $k[[D]]$.

After faithfully flat trivialization of a finite torsor, its function
module is the dual regular module of the truncated polynomial ring.
The pairing that extracts the coefficient of $D^{p^{2l}-1}$ from a
product identifies that dual module with a single regular nilpotent
block.  Consequently $\ker D^e$ on $R_l$ has rank $e$ over $L$ for
$1\le e\le p^{2l}$.  These ranks descend under faithful flatness.
For $q=p^i$, the formal-group coproduct satisfies
\[
 \Delta(D^q)\in(D^q\otimes1,1\otimes D^q).
\]
Hence $\ker D^q$ is a subalgebra.  On a field factor it is an
intermediate field of purely inseparable degree $q$.

The element $t$ is a $p$-basis of $K$ and remains a $p$-basis under
separable extension.  An intermediate field of degree $p^i$ in a
one-generator purely inseparable extension is contained in
$L^{1/p^i}$: every element has purely inseparable degree at most
$p^i$.  Equality follows from the degrees.  This proves the kernel
formula on field factors, and hence on $L$.  To justify the last
passage without requiring a chosen global field factor, note that
$L$ is an ind-\'{e}tale algebra over a field and is absolutely flat.
The finite free module identities in question can be checked at its
residue fields, which are separable algebraic extensions of $K$.
In a larger nilpotent block, the old kernel is contained in the image
of $D$.  This proves surjectivity on the union and nonvanishing of
the top operator on every finite kernel.

By Lemma~\ref{lem:etale-marking}, a norm unit rescales the chosen
\'{e}tale generator by that unit or its inverse.  It therefore
conjugates the translation group by the corresponding scalar.
Lemma~\ref{lem:lift-torsors} checks directly that the central
order-two unit acts by $[-1]$ on the named translations.  On the
tangent of the Cartier dual it therefore has character $-1$.
If $D_0$ is any parameter and $\tau$ denotes that order-two
action, replace it by $(D_0-\tau(D_0))/2$.  Its linear coefficient is
unchanged, so it is a parameter and satisfies $\tau(D)=-D$.
Continuity of the marking action gives continuity on every finite jet
of this parameter algebra.

For the Frobenius identities, use the Honda model over $\Fp$.
Its Frobenius $\Phi$ satisfies $\Phi^2=[p]$.  Since
$V\Phi=[p]=\Phi^2$ on the divisible group and $\Phi$ is faithfully
flat, cancellation gives $V=\Phi$.  Cartier duality exchanges these
operators.  For the indicated even $i$, their $i$th iterates are
$[p^{i/2}]$.  If $A_l=\mathcal O(\Gamma_2[p^l])$ and
$\langle b,\xi\rangle=\xi(b)$ is its pairing with distributions,
Cartier-dual naturality gives, because $q$ fixes $k$,
\[
 \langle b^q,D\rangle
 =\langle b,(V^i)^*D\rangle
 =\langle b,D^q\rangle.
\]
For the last equality, an $\Fp$-coordinate $z$ on the dual has
$(V^i)^*z=z^q$; writing $D=\sum c_nz^n$ over $k$ gives
$(V^i)^*D=D^q$.  If a torsor coaction is
$\rho(f)=\sum_j f_j\otimes b_j$, it follows that
\[
 D(f^q)=\sum_j f_j^q\langle b_j^q,D\rangle
       =\left(\sum_j f_j\langle b_j,D^q\rangle\right)^q
       =(D^qf)^q.
\]
The coefficients $f_j$ retain their $q$th powers; no splitting of the
torsor is used.  Apply this identity to $af$ with $a^q\in L$ and use
$L$-linearity of $D$ to obtain the second identity in
\eqref{eq:distribution-frobenius}.

On a fixed root level only finitely many powers of $D$ act.  Their
matrices in the fractional-power basis of the $p$-basis have finitely
many coefficients in $L$, hence are defined at one finite separable
stage.  They are continuous maps between finite-dimensional spaces
over complete local fields.  For a finite target stage, choose a
$D$-preimage of each member of a finite basis over that stage's
separable coefficient algebra.  Put these finitely many preimages
in one larger stage and extend linearly.  This is a continuous
additive section.  At a finite root stage the intersection with
$\ker D=L$ is its separable stage, with the induced topology, by
separability versus pure inseparability.  The coefficient exactness
criterion above proves the final assertion, including parameters.
\end{proof}

\begin{proposition}[An invariant differential]
\label{prop:invariant-differential}
There is an $H$-invariant basis
$\eta\in\Omega^1_{L/k}$ defined, together with its inverse frame,
at one finite separable marking stage.
\end{proposition}

\begin{proof}
By Lemma~\ref{lem:finite-p-basis}, ordinary differentials of $A$ are freely
generated by $da,dt$ and agree with continuous differentials.
The Kodaira--Spencer
homomorphism for the universal $p$-divisible group is the functorial
isomorphism
\begin{equation}
 \label{eq:KS-decompletion}
 \Hom_A(\Lie\mathcal G,\omega_{\mathcal G^\vee})
 \xrightarrow{\ \sim\ }\Omega^1_{A/k}.
\end{equation}
For a $p$-divisible group $G$ over any ring $B$ on which $p$ is
nilpotent, isomorphism classes of lifts across $B=B'/I$ with $I^2=0$
form a torsor under
$\Hom_B(\omega_{G^\vee},I\otimes_B\Lie G)$.  The associated functorial
Kodaira--Spencer map
$\Hom_B(\Lie G,\omega_{G^\vee})\to\Omega^1_{B/\mathbb Z}$ controls
changes of a lifting ring map
\cite[Theorem~5.1 and Equation~(5.1), p.~226]{Lau10}.
The isomorphism in \eqref{eq:KS-decompletion} uses the separate
universal-deformation case
\cite[paragraph following Equation~(5.2), p.~227]{Lau10}:
$A=k[[a,t]]$ is a complete local Noetherian $k$-algebra with residue
field $k$, the deformation is universal, and $k$ is perfect.
In this case Lau identifies completed relative differentials with
ordinary absolute differentials.  Since every absolute derivation kills
the perfect field $k$, the latter are $\Omega^1_{A/k}$, in agreement
with the finite $p$-basis calculation above.  Universality makes the
closed-point map bijective; both modules are free of rank two, so the
map is an isomorphism over $A$.

Over $L$ the two markings give
\[
 0\longrightarrow\Gamma_{2,L}
 \longrightarrow\mathcal G_L
 \longrightarrow(\Qp/\Zp)_L\longrightarrow0.
\]
Thus $\Lie\mathcal G_L=\Lie\Gamma_{2,L}$ has an $H$-fixed frame.
Dualizing gives the exact sequence of invariant-form modules
\begin{equation}
 \label{eq:dual-hodge-marked-sequence}
 0\longrightarrow\omega_{\Gamma_2^\vee,L}
 \longrightarrow\omega_{\mathcal G^\vee,L}
 \longrightarrow\omega_{\mu_{p^\infty},L}
 \longrightarrow0.
\end{equation}
Both end lines have $H$-fixed frames.  Exactness may be checked after
the faithfully flat extension to perfect coefficients, where the
marked connected--\'{e}tale extension splits, and then descended.
The determinant of the middle module is therefore equivariantly
trivial even when the extension over $L$ is nonsplit.  Taking
determinants in \eqref{eq:KS-decompletion} shows that
$\det(\Omega^1_{A/k}\otimes_A L)$ has an $H$-fixed basis.

The conormal line of $a=0$ also has an $H$-fixed frame on this marking
cover.  Explicitly, comparison of the coefficient of $T^p$ under a
coordinate change with leading coefficient $c$ gives
$a'=c^{1-p}a$ in the corresponding coordinate convention.  Thus
$a(dT)^{p-1}$ is the first Hasse section, and its conormal has the
inverse periodicity twist.  The connected marking frames $dT$ and
trivializes that twist.  More explicitly, if the fixed cotangent
frame is $e=u\,dT$, then $T'=cT+\cdots$ gives $u'=u/c$ and
$a'=c^{1-p}a$.  Hence $[a'](u')^{1-p}=[a]u^{1-p}$ is an
invariant conormal frame.  Since $L/K$ is ind-\'{e}tale, the conormal
sequence is
\[
 0\longrightarrow L\,da
 \longrightarrow\Omega^1_{A/k}\otimes_A L
 \longrightarrow\Omega^1_{L/k}\longrightarrow0.
\]
Taking the determinant quotient produces the desired $H$-fixed
basis $\eta$.  This construction uses differentials over $A$ and
$L$, not differentials of a perfection.  Writing $\eta=f\,dt$, both
$f$ and $f^{-1}$ belong to a common finite marking stage; enlarging
that stage if necessary makes $\eta$ a basis on every field factor.
\end{proof}

Theorem~\ref{thm:completed-calculation} gives, before decompletion,
\begin{equation}
 \label{eq:L-H-invariants}
 L^H=k.
\end{equation}
Write $\mathrm{Fr}(f)=f^p$ for coefficient Frobenius and
$\mathcal O_B$ for the product of the valuation rings of the field
factors of a finite marking stage $B$.
Let $\operatorname{Car}$ be intrinsic Cartier on differentials
\citep[Chapter~II, Section~6]{Cartier58}, and set
\begin{equation}
 \label{eq:normalized-Cartier}
 C_0(f)=\frac{\operatorname{Car}(f\eta)}{\eta}.
\end{equation}
It is $\Fp$-linear, satisfies $C_0(a^pf)=aC_0(f)$ for $a,f\in L$,
and commutes with $H$.

\begin{lemma}[Cartier traces and residue duality]
\label{lem:Cartier-traces}
For every $i\ge1$, with $q=p^i$, there is $c_i\in k^\times$ such that
\begin{equation}
 \label{eq:Cartier-distribution-trace}
 D^{q-1}(f^{1/q})=c_i C_0^i(f)\qquad(f\in L).
\end{equation}
In particular, $C_0(f^p)=0$.  If $B=L^U$ contains $\eta$ as a
basis, then
\begin{equation}
 \label{eq:Cartier-exact-sequence}
 0\longrightarrow B\xrightarrow{\mathrm{Fr}}B
 \xrightarrow{d/\eta}B\xrightarrow{C_0}B\longrightarrow0
\end{equation}
is exact and its two short exact factors have continuous additive
sections onto their images.

Write $B=\prod_j\kappa_j((s_j))$ and
$\eta=h_j(s_j)\,ds_j/s_j$ on its factors.  Define
\[
 M=\{f:\operatorname{ord}_{s_j}(f_jh_j)>0\text{ for all }j\},
 \qquad
 M^0=\{f:\operatorname{ord}_{s_j}(f_jh_j)\ge0\text{ for all }j\}.
\]
These are compact open $H$- and $C_0$-stable lattices, and $M^0/M$
is finite.  For the pairing
\begin{equation}
 \label{eq:residue-pairing}
 \langle f,g\rangle_B
 =\sum_j\operatorname{Tr}_{\kappa_j/\Fp}
             \operatorname{Res}_{s_j}(f_jg_j\eta)
\end{equation}
there are topological $H$-equivariant identifications
\begin{equation}
 \label{eq:residue-lattice-duals}
 M^\vee=B/\mathcal O_B,\qquad
 (B/M)^\vee=\mathcal O_B,\qquad
 \langle C_0f,g\rangle_B=\langle f,g^p\rangle_B.
\end{equation}
\end{lemma}

\begin{proof}
The functional $I_i(f)=D^{q-1}(f^{1/q})$ takes values in $L$ by
\eqref{eq:distribution-kernels}; it is nonzero on every field factor
by the nilpotent-block calculation.  It satisfies
$I_i(a^qf)=aI_i(f)$ and is $H$-equivariant.  The same semilinearity
and equivariance hold for $C_0^i$.  The Cartier pairing gives an
isomorphism
\begin{equation}
 \label{eq:Cartier-perfect-pairing}
 \mathrm{Fr}_*^iL\longrightarrow
 \Hom_L(\mathrm{Fr}_*^iL,L),\qquad
 b\longmapsto\bigl[f\longmapsto C_0^i(bf)\bigr].
\end{equation}
Here $\mathrm{Fr}_*^iL$ has the scalar action $a\cdot f=a^qf$.
Both modules are free of rank $q$ in the common $p$-basis.  On a
field factor the pairing is nondegenerate: for $b\ne0$, choose $z$
with $C_0^i(z)\ne0$ and evaluate at $b^{-1}z$.  This proves
\eqref{eq:Cartier-perfect-pairing}, also for finite products and the
ind-\'{e}tale union.  Thus a unique multiplier $b_i$ satisfies
$I_i(f)=C_0^i(b_if)$.  It is a unit by nonvanishing on each factor.
Equivariance and uniqueness give $b_i\in L^H=k$ by
\eqref{eq:L-H-invariants}.  Taking $c_i=b_i^{1/q}$ proves
\eqref{eq:Cartier-distribution-trace}.

This argument works for $i=1$ without the even-iterate restriction
in \eqref{eq:distribution-frobenius}.  Since $D$ kills $L$, it gives
$c_1C_0(f^p)=D^{p-1}f=0$.  In particular $C_0\mathrm{Fr}=0$ is a
consequence of the invariant torsor and differential; it was not
assumed for an arbitrary choice of $\eta$.

On a Laurent field the Cartier formula is
\[
 \operatorname{Car}\left(\sum_n a_ns^n\frac{ds}{s}\right)
 =\sum_n a_{pn}^{1/p}s^n\frac{ds}{s}.
\]
It proves continuity and surjectivity.  Its kernel consists of
differentials whose logarithmic coefficients at indices divisible
by $p$ vanish; a continuous additive primitive is
$\sum_{p\nmid n}(a_n/n)s^n$.  The kernel of $d$ is $B^p$, and its
root map is continuous.  A continuous section of Cartier sends
$\sum a_ns^n ds/s$ to $\sum a_n^ps^{pn}ds/s$.
Multiplication or division by $\eta$ gives the sections and
exactness in \eqref{eq:Cartier-exact-sequence}.

The logarithmic order of a differential is independent of the local
uniformizer, since $ds'/s'$ is a unit multiple of $ds/s$.
Invariance of $\eta$ therefore makes $M$ and $M^0$ invariant under
$H$, including permutations of the field factors.  The Cartier
formula preserves positive and nonnegative logarithmic order.
Furthermore $M^0/M\cong\prod_j\kappa_j$ as additive groups.

Residue is invariant under coordinate change, and the sum of traces
is invariant under residue-field automorphisms and permutations.
Thus \eqref{eq:residue-pairing} is $H$-invariant.  A continuous
functional on $M$ uses finitely many Laurent coefficients; it is
represented by an element of $B/\mathcal O_B$.  Conversely every
negative principal part is detected by pairing against a suitable
positive logarithmic coefficient.  This proves the first dual
identification.  An arbitrary functional on the discrete space
$B/M$ is a sequence of coefficient functionals, represented by a
power series in $\mathcal O_B$, which proves the second.  The classical
residue adjointness includes a Frobenius twist
\citep[Section~10, Proposition~9]{Serre58Topology}.  Here the local
normalization gives
\[
 \operatorname{Res}\bigl(g\operatorname{Car}(f\eta)\bigr)
 =\operatorname{Res}\bigl(\operatorname{Car}(g^pf\eta)\bigr).
\]
Taking the finite-field trace removes the $p$th-root operation on
the residue and proves the last identity in
\eqref{eq:residue-lattice-duals}.
\end{proof}

\begin{lemma}[Good finite stages]
\label{lem:good-stages}
There is a cofinal collection of open normal uniform subgroups
$U\subset J$ for which $B=L^U$ has all the following properties:
\begin{enumerate}
\item $\eta$ is a basis over $B$ and $H^a_{\mathrm{cts}}(H,B^\flat)$
is finite for every $a$;
\item there is $\theta_0\in B^{1/p}$ with $D\theta_0=1$;
\item for every $\sigma\in U$,
\begin{equation}
 \label{eq:good-stage-jets}
 \sigma(D)-D\in D^{p+2}k[[D]].
\end{equation}
\end{enumerate}
For these stages $D$ restricts to a continuous $H$-equivariant map
$B^{1/p}\to B^{1/p}$.
\end{lemma}

\begin{proof}
In a regular nilpotent block, a preimage $\theta_0$ of $1$ is killed
by $D^2$ and hence by $D^p$.  The kernel formula puts it in
$L^{1/p}$.  Its $p$th power and the finite coefficients of $\eta$
and its inverse belong to one finite separable stage.  The action on
finite distribution jets is continuous, so the kernel of its action
modulo $D^{p+2}$ is open.  Intersect these finite-data stabilizers,
take an open normal core, and pass to a still smaller uniform subgroup
in the cofinal collection of
Theorem~\ref{thm:completed-calculation}.  This gives all three
conditions, cofinally.

If $b\in B^{1/p}$, then $D^pb=0$.  For $\sigma\in U$,
\eqref{eq:good-stage-jets} gives $\sigma(Db)=Db$; moreover $Db$ is
still killed by $D^p$.  Uniqueness of purely inseparable roots
identifies $(L^{1/p})^U=B^{1/p}$, so $D$ preserves this stage.
Continuity follows from Lemma~\ref{lem:distribution-operator} and
the induced valuation topology.
\end{proof}

\subsection{Finiteness at a separable stage}

For any of our valued algebras $E$, write
$E^{++}=\{f:v(f)>0\}$, where positivity is required on every field
factor, and write $\mathcal O_E$ for the elements of nonnegative
valuation.  These are the ordinary valuation lattices.  The lattice
$M$ in Lemma~\ref{lem:Cartier-traces}, by contrast, measures the
logarithmic order of $f\eta$.

The completed calculation first gives finite cohomology for
$B^\flat/\mathcal O_{B^\flat}$.  Residue duality will turn this
into a finite stable Cartier image on the compact cohomology of $M$.
We record that completed-coefficient input and a compact-operator
lemma before using them to prove finiteness at $B$.

\begin{lemma}[Completed positive coefficients]
\label{lem:completed-positive}
At every good stage $B$, the complex
$\Ccts(H,(B^\flat)^{++})$ is acyclic.  Moreover,
$H^a(H,\mathcal O_{B^\flat})$ and
$H^a(H,B^\flat/\mathcal O_{B^\flat})$ are finite for every $a$.
\end{lemma}

\begin{proof}
Write $\partial$ for the cochain differential and $\mathrm{Fr}$ for
coefficient Frobenius.  A continuous positive-valued cochain on
the compact space $H^r$ has valuations bounded below by some
$\epsilon>0$.  Its iterates under $\mathrm{Fr}$ therefore tend
uniformly to zero.  In positive degree, apply
Lemma~\ref{lem:strictness} to $\Ccts(H,B^\flat)$, whose cohomology is
finite by Theorem~\ref{thm:completed-calculation}.  A sufficiently
large Frobenius iterate of a positive cocycle has a primitive with
positive values.  Frobenius is a homeomorphism of $B^\flat$ and of
its positive lattice; its inverse carries that primitive back to a
positive primitive of the original cocycle.  In degree zero the
finite invariant group is discrete in the cycle topology.  A
sufficiently large Frobenius iterate of a positive invariant is
therefore zero, and injectivity of Frobenius makes the invariant zero.

The residue algebra
$\mathcal O_{B^\flat}/(B^\flat)^{++}$ is a finite product of finite
fields.  Its cohomology is finite by
Lemma~\ref{lem:finite-resolution}.  The valuation quotients have
continuous sections, so the long exact sequences for the positive,
integral, and full coefficients prove both finiteness assertions.
\end{proof}

We isolate the compact algebra used in the countability argument.

\begin{lemma}[A compact operator]
\label{lem:compact-operator}
Let $N$ be a compact Hausdorff $\Fp$-vector space, let $C:N\to N$
be continuous and $\Fp$-linear, and suppose that
\[
 I=\bigcap_{n\ge0}C^nN
 \quad\text{and}\quad N/CN
\]
are finite.  Then $\ker C$ is finite, and every subgroup of $N$
whose elements have finite forward $C$-orbits is finite.  If
$f:N\to X$ is a homomorphism with finite kernel, then
\begin{equation}
 \label{eq:stable-image-finite-fiber}
 \bigcap_{n\ge0}f(C^nN)=f(I).
\end{equation}
No topology on $X$ is required in this last assertion.
\end{lemma}

\begin{proof}
Compactness of the nested sets of preimages of an element of $I$
gives $CI=I$.  On $Q=N/I$, the compact images $C^nQ$ have
intersection zero.  They eventually lie in any specified
neighborhood of zero: otherwise their intersections with the
closed complement of that neighborhood would have the finite
intersection property.  Thus $C$ is uniformly topologically
nilpotent on $Q$.  The formula
\[
 \left(\sum_{n\ge0}a_nz^n\right)x
   =\sum_{n\ge0}a_nC^nx
\]
defines a continuous $\Fp[[z]]$-module structure on $Q$, with $z=C$.
The series converges uniformly, and this also proves continuity.

Choose lifts $x_1,\ldots,x_r$ of a basis of the finite space $Q/CQ$.
Successively express a vector modulo $CQ$, then its remainder
modulo $C^2Q$, and so on.  The remainder tends uniformly to zero,
so every vector is an $\Fp[[z]]$-linear combination of the $x_j$.
Hence $Q$ is a finitely generated module over the discrete valuation
ring $\Fp[[z]]$.  Its structure theorem makes its $z$-power torsion
finite.  In particular $\ker(C:Q\to Q)$ is finite, and the finite
kernel of $N\to Q$ then makes $\ker(C:N\to N)$ finite.

If an element has finite forward orbit, its image in $Q$ tends to
zero through a finite set and is eventually zero.  Such images
belong to the finite $z$-power torsion of $Q$.  The fibers of
$N\to Q$ are finite, proving the orbit assertion.

For the last claim, fix $x$ in the left side of
\eqref{eq:stable-image-finite-fiber}.  The sets
$f^{-1}(x)\cap C^nN$ are nonempty nested compact sets: each is a
closed subset of the finite fiber $f^{-1}(x)$.  Their intersection
is nonempty and lies in $I$.  The reverse inclusion is immediate.
\end{proof}

\begin{proposition}[Finite-stage finiteness]
\label{prop:finite-stage-finiteness}
For every good $B$, every finite root stage $B^{1/p^e}$, and
every $H$-stable valuation-order lattice in such a stage,
continuous $H$-cohomology is finite in every degree.
\end{proposition}

\begin{proof}
We first prove countability at the separable stage $B$.  The finite
cohomological range lets us argue at a largest degree where
countability could fail.  At that degree, Cartier will force
Frobenius to have countable image; an explicit distribution lift will
show that its kernel is countable as well.  Once countability is
known in every degree, compact duality will give finiteness.

Fix a good stage $B$ and set
\begin{equation}
 \label{eq:finite-stage-notation}
 N_a=H^a(H,M),\qquad X_a=H^a(H,B),\qquad
 f_a:N_a\longrightarrow X_a.
\end{equation}
The spaces $N_a$ are compact Hausdorff by
Lemma~\ref{lem:compact-duality}; no separatedness is yet asserted
for $X_a$.  Cartier acts on both and commutes with $f_a$.
Since $B/M$ is countable discrete, its cochain groups are countable.
The long exact sequence therefore gives
\begin{equation}
 \label{eq:countable-defects}
 \ker f_a\ \text{and}\ \coker f_a\ \text{are countable}.
\end{equation}

\smallskip\noindent\emph{Step 1: the compact Cartier model.}
For every degree $a$, put
\[
 I_a=\bigcap_{n\ge0}C_0^nN_a.
\]
We claim that $I_a$ is finite.  The completed calculation supplies
finite cohomology of the discrete principal-part quotient.

Compact duality and the residue pairing identify
\[
 N_a^\vee=H^{8-a}(H,B/\mathcal O_B),
\]
and transpose $C_0$ to Frobenius.  There is an identification of
discrete $H$-modules
\begin{equation}
 \label{eq:principal-parts-perfection}
 \colim_{\mathrm{Fr}} B/\mathcal O_B
 =B^{\mathrm{perf}}/\mathcal O_{B^{\mathrm{perf}}}
 =B^\flat/\mathcal O_{B^\flat}.
\end{equation}
At the $e$th copy of the left side the first map sends $f$ to
$f^{1/p^e}$; its compatibility is precisely the Frobenius
transition.  The second equality follows by density: every
completed element differs from an element of the root-stage union
by an integral element, and integrality on that union is the
induced valuation condition.

Continuous cochains into discrete modules commute with filtered
colimits, since they have finite image.  Consequently
$\colim_{\mathrm{Fr}}N_a^\vee$ is finite by
Lemma~\ref{lem:completed-positive}.  Dualizing makes
$\varprojlim_{C_0}N_a$ finite.  Its projection onto its zeroth
coordinate is exactly $I_a$.  To see surjectivity onto $I_a$,
the finite strings of compatible predecessors of any element of
$I_a$ form nonempty compact sets; compactness supplies an infinite
compatible string.  Thus $I_a$ is finite.

\smallskip\noindent\emph{Step 2: Cartier at a largest uncountable degree.}
Suppose that some $X_a$ is uncountable and
choose the largest such $a$; such an $a$
exists because cohomology vanishes above eight.  Split
\eqref{eq:Cartier-exact-sequence} into the two short exact
sequences with intermediate module $V=\im(d/\eta)$.
They are exact on cochains.  The first sequence shows that
$H^{a+1}(H,V)$ is countable, since it lies between quotients
or subgroups of the countable groups $X_{a+1}$ and $X_{a+2}$.
The second sequence embeds $X_a/C_0X_a$ in
$H^{a+1}(H,V)$, so this quotient is countable.

Put $W=\ker f_a$, $A_a=\im f_a$, and $E_a=\coker f_a$.
The exact sequences for $C_0$ on
$0\to W\to N_a\to A_a\to0$ and
$0\to A_a\to X_a\to E_a\to0$ show that
$N_a/C_0N_a$ is countable: its successive possible contributions
are $W/C_0W$, $E_a[C_0]$, and $X_a/C_0X_a$, all countable by
\eqref{eq:countable-defects} and the countability of $X_a/C_0X_a$ just proved.
This quotient is compact Hausdorff.  A countable compact
Hausdorff group is finite by Baire's theorem: if no singleton
were open, its countable cover by singletons would be meagre.
Lemma~\ref{lem:compact-operator} now applies to $N_a$, since
Step~1 made its stable Cartier image $I_a$ finite.

Every element of $W$ has finite forward Cartier orbit.  In degree
zero $W=0$.  In positive degree represent it by the connecting
image of a cocycle with values in $B/M$.  That cochain has finite
image.  In the logarithmic Laurent expansion of any of its
finitely many values, only finitely many negative indices occur.
Repeated Cartier eventually removes all of them, leaving values
in the finite module $M^0/M$.  The resulting cohomology classes
belong to the image of the finite group
$H^{a-1}(H,M^0/M)$.  Cartier preserves this image.  Hence the
forward orbit is finite, and Lemma~\ref{lem:compact-operator}
makes $W$ finite.

The same lemma gives a finite kernel of $C_0$ on $N_a$.
The kernel on $A_a$ is then finite, by the exact sequence with
the now finite group $W/C_0W$.  The kernel on $X_a$ maps into
the countable group $E_a[C_0]$, so it is countable.
Finally, the finite kernel of $f_a$ allows us to apply
\eqref{eq:stable-image-finite-fiber}, giving
\begin{equation}
 \label{eq:finite-stable-X-image}
 \bigcap_{n\ge0}f_a(C_0^nN_a)=f_a(I_a),
\end{equation}
a finite group.  This will control the Frobenius-boundary presentations
in the next step.

\smallskip\noindent\emph{Step 3: the Frobenius kernel at that degree.}
We now prove that $\ker(\mathrm{Fr}:X_a\to X_a)$ is countable.
A Frobenius boundary has a presentation in $B^{1/p}$; applying $D$
turns it into a cocycle.  Its principal part has only countably many
possibilities.  The remaining positive-valued presentations will
all map into the fixed finite group in
\eqref{eq:finite-stable-X-image}.

In degree zero the kernel is zero by injectivity of coefficient
Frobenius.  For $a>0$, consider the additive presentation group
\[
 \mathcal P_a=
 \{b\in C^{a-1}_{\mathrm{cts}}(H,B^{1/p}):
                  \partial b\in C^a_{\mathrm{cts}}(H,B)\}.
\]
The map $b\mapsto[\partial b]$ takes values in the Frobenius kernel,
since $(\partial b)^p=\partial(b^p)$ and $b^p$ is $B$-valued.
It surjects onto that kernel:
if $\mathrm{Fr}(\alpha)=\partial c$ for a $B$-valued cocycle
$\alpha$, take $b=c^{1/p}$.  For any presentation put $w=Db$.
The good-stage property puts $w$ in
$C^{a-1}_{\mathrm{cts}}(H,B^{1/p})$, and $\partial w=0$
because $D$ kills $B$ and commutes with $H$.

The quotient $B^{1/p}/(B^{1/p})^{++}$ is countable discrete.
Thus the map
\[
 \mathcal P_a\longrightarrow
 C^{a-1}_{\mathrm{cts}}
       (H,B^{1/p}/(B^{1/p})^{++}),\qquad b\longmapsto Db
\]
has countable image.  Its kernel $\mathcal P_a^+$, consisting
of presentations with positive $w$, has countable index.
It suffices to show that $\mathcal P_a^+$ has finite image in
$X_a$.

Fix $b\in\mathcal P_a^+$.  If $w=0$, then $b$ is $B$-valued
by \eqref{eq:distribution-kernels}, and $[\partial b]=0$.
Otherwise compactness gives a uniform lower bound
$v(w)\ge\epsilon>0$ on all field factors and cochain values.
Let $i$ be positive, even, and divisible by $[k:\Fp]$, and put
$q=p^i$.  With the good-stage element $\theta_0$ define
\begin{equation}
 \label{eq:Frobenius-presentation-lift}
 b'=\theta_0^{1/q}w
 \in C^{a-1}_{\mathrm{cts}}(H,B^{1/(pq)}).
\end{equation}
Since $w\in L^{1/q}$, the Frobenius identities give
$D^qb'=w$.  Therefore $\partial b'$ is killed by $D^q$,
so it belongs to $L^{1/q}$.  It is also fixed by $U$, hence
is $B^{1/q}$-valued.

Similarly $e=D^{q-1}b'-b$ is killed by $D$.  It is fixed by
$U$ as well.  Indeed \eqref{eq:good-stage-jets} gives
\[
 \sigma(D)=D(1+h),\quad h\in D^{p+1}k[[D]],
 \qquad
 \sigma(D^{q-1})-D^{q-1}\in D^{q+p}k[[D]].
\]
The last ideal kills $b'$, since $D^{q+p}b'=D^pw=0$.
Thus $e$ is $B$-valued.  Every operation here is continuous
at a finite marking and root stage by
Lemma~\ref{lem:distribution-operator}; once its values have
been shown to lie in $B$ or $B^{1/q}$, the induced topology
makes it a continuous cochain there.  The auxiliary larger
stage may depend on $i$.

Set $\alpha_i=(\partial b')^q$, a $B$-valued cocycle.
The equation $\partial e=D^{q-1}\partial b'-\partial b$ and
Lemma~\ref{lem:Cartier-traces} give
\begin{equation}
 \label{eq:Frobenius-kernel-Cartier}
 [\partial b]=c_i C_0^i[\alpha_i]\quad\text{in }X_a.
\end{equation}
Normalize valuations by their restriction to $K$ and take their
minimum over the factors.  The $H$-action preserves these
valuations, so the cochain differential does not lower the
minimum.  Formula~\eqref{eq:Frobenius-presentation-lift} yields
\[
 v(\alpha_i)\ge q\epsilon+v(\theta_0).
\]
For every sufficiently large admissible $i$ this exceeds the
fixed bounds defining $M$, and $\alpha_i$ is $M$-valued.
Since $c_i\in k^\times$,
$c_iC_0^i(\alpha_i)=C_0^i(c_i^q\alpha_i)$, with
$c_i^q\alpha_i$ still in $M$.
It follows that $[\partial b]\in f_a(C_0^iN_a)$ for a
cofinal sequence of $i$.  These images are nested, so
\eqref{eq:finite-stable-X-image} puts every such class in the
fixed finite group $f_a(I_a)$.  This proves that the positive
presentation subgroup has finite image.  Its countable index
therefore makes the Frobenius kernel countable, as required.

\smallskip\noindent\emph{Step 4: from countability to finiteness.}
At the degree chosen in Step~2, $C_0\mathrm{Fr}=0$ puts the
Frobenius image inside the countable group
$\ker(C_0:X_a\to X_a)$.  Step~3 makes its kernel countable.
Thus $X_a$ is countable, contrary to its choice.  Every $X_a$
is therefore countable.

For every degree $a$, \eqref{eq:countable-defects} now makes
$N_a$ countable.
It is compact Hausdorff, hence finite.  Any $H$-stable
valuation-order lattice is commensurable with $M$, with
finite quotients after intersecting the two lattices.
Long exact sequences and finite-coefficient cohomology
therefore give finite cohomology for every such lattice,
including $\mathcal O_B$.
The second residue duality in
\eqref{eq:residue-lattice-duals}, followed by
Lemma~\ref{lem:compact-duality}, makes
$H^a(H,B/M)$ finite.  The long exact sequence for
$0\to M\to B\to B/M\to0$ now makes each $X_a$ finite.

Finally $p^e$th power is an additive $H$-equivariant
homeomorphism $B^{1/p^e}\to B$.  It takes valuation-order
lattices to valuation-order lattices and transports all
the asserted finiteness statements.
\end{proof}

\subsection{The natural decompletion map}

\begin{theorem}[Decompletion, with profinite parameters]
\label{thm:decompletion}
For every compact profinite space $T$, the natural coefficient
inclusions induce a quasi-isomorphism
\begin{equation}
 \label{eq:natural-decompletion}
 \Ccts(H,R;T)
 =\colim_{U,e}\Ccts(H,(L^U)^{1/p^e};T)
 \longrightarrow
 \colim_U\Ccts(H,(L^U)^\flat;T).
\end{equation}
It is natural in $T$, in marking-stage inclusions, and for the
$P$- and $J$-actions.  In particular,
\begin{equation}
 \label{eq:perfect-H-cohomology}
 H^a\Ccts(H,R;T)=
 \begin{cases}
 C_{\mathrm{cts}}(T,k),&a=0,\\
 C_{\mathrm{cts}}(T,\ell),&a=3,\\
 0,&a\ne0,3.
 \end{cases}
\end{equation}
The two lines are discrete.  The residual $P/H$-action on them
is trivial, $J$ acts trivially on $k$, and
$\mathcal O_F^\times\subset J$ acts on $\ell$ by $\delta$.
\end{theorem}

\begin{proof}
First fix a good stage $B$ and take $T$ to be a point.
Proposition~\ref{prop:finite-stage-finiteness} gives finite
cohomology of $B^{++}$.  A positive cocycle tends uniformly to
zero under successive Frobenius powers.  By
Lemma~\ref{lem:strictness}, its sufficiently large iterates are
boundaries in the positive complex.  Thus Frobenius acts
nilpotently on each finite group $H^a(H,B^{++})$.
The same conclusion in degree zero follows from discreteness
of the finite invariant group and injectivity of Frobenius.
Transporting root-stage inclusions by their power
homeomorphisms identifies their action on cohomology with
Frobenius.  It follows that
\begin{equation}
 \label{eq:positive-root-acyclicity}
 \colim_e\Ccts(H,(B^{1/p^e})^{++})
 \quad\text{is acyclic}.
\end{equation}
The completed positive complex is acyclic by
Lemma~\ref{lem:completed-positive}.

Density gives an isomorphism of discrete $H$-modules
\begin{equation}
 \label{eq:positive-quotient-comparison}
 B^{\mathrm{perf}}/(B^{\mathrm{perf}})^{++}
 \xrightarrow{\ \sim\ }
 B^\flat/(B^\flat)^{++}.
\end{equation}
Surjectivity follows by approximating a completed element
with error of strictly positive valuation; injectivity is
the induced valuation condition.  A continuous map from a
compact space into either quotient has finite image.
Its finitely many values lift to one root stage on the left
and to the completed algebra on the right.  Hence the
coefficient sequences with positive kernels are exact on
the prescribed cochains.  The acyclicity of their kernels
and \eqref{eq:positive-quotient-comparison} prove the
fixed-$B$ comparison.

Now allow arbitrary compact profinite $T$.
All finite-root positive complexes have finite cohomology,
so Lemma~\ref{lem:profinite-parameters} upgrades
\eqref{eq:positive-root-acyclicity} to the corresponding
parameterized assertion.  The same lemma upgrades completed
positive acyclicity.  The finite-image lifting argument for
\eqref{eq:positive-quotient-comparison} applies to
$T\times H^r$ without any countability assumption on $T$.
It proves the fixed-$B$ comparison with these parameters.

Taking the filtered colimit over the cofinal family of
good stages proves \eqref{eq:natural-decompletion}.
The comparison itself is the coefficient inclusion on the
full diagram of marking and root stages.  The differential
$\eta$, the parameter $D$, and the good subfamily were used
only to prove it is a quasi-isomorphism.  Consequently its
naturality for the full $P$- and $J$-actions is unaffected by
these choices or by conjugating to another cofinal family.
Finally apply Theorem~\ref{thm:completed-calculation} to
obtain \eqref{eq:perfect-H-cohomology} and all the stated
actions.
\end{proof}

To pass from $H$ to the norm quotients, we spell out the
continuous extension argument.  This also specifies why
the parameterized comparison is the needed one.

\begin{lemma}[Continuous extension descent]
\label{lem:continuous-extension}
Let $1\to N\to G\to Q\to1$ be an extension of compact profinite groups
whose quotient map has a continuous section as a map of
spaces.  Let $M$ be a complete continuous coefficient
module with invariant neighborhoods, or a filtered union
of complete $G$-stable stages with the convention of
Definition~\ref{def:stage-cochains}.

If continuous $N$-cohomology satisfies
\[
 H^t\Ccts(N,M;T)=C_{\mathrm{cts}}(T,V_t)
\]
for all compact profinite $T$, naturally in $T$ and for
the quotient action, with $V_t$ discrete, there is the
natural first-quadrant spectral sequence
\begin{equation}
 \label{eq:continuous-HS}
 H^s_{\mathrm{cts}}(Q,V_t)
 \Longrightarrow H^{s+t}\Ccts(G,M).
\end{equation}
The same construction includes an additional profinite
parameter.  More generally, a coefficient map inducing
quasi-isomorphisms on $N$-cochains with all such parameters
induces a quasi-isomorphism on $G$-cochains.
\end{lemma}

\begin{proof}
We give a bar-resolution construction that works on
each fixed coefficient stage.  Put
\[
 I^j(M)=C_{\mathrm{cts}}(G^{j+1},M),\qquad
 (g f)(g_0,\ldots,g_j)
   =g\,f(g^{-1}g_0,\ldots,g^{-1}g_j),
\]
with the alternating omission differential.  A section
$s:Q\to G$, chosen with $s(1)=1$, gives an $N$-equivariant
continuous retraction
$r(g)=g\,s(\overline g)^{-1}$ from $G$ to $N$.
Restriction to $N^{j+1}$ and pullback along $r$ compare
$I^\bullet(M)^N$ with homogeneous $N$-cochains.
They are mutually inverse up to a continuous cochain
homotopy: the usual prism inserts, one vertex at a time,
the vertices $r(g_i)$ in place of $g_i$.
In degree $j$ this is the finite alternating sum obtained
by evaluating at
\[
 (g_0,\ldots,g_i,r(g_i),\ldots,r(g_{j-1})),
 \qquad 0\le i<j.
\]
Cancellation of adjacent omissions gives the homotopy
identity.  It preserves a fixed coefficient stage and
commutes with a profinite parameter.  Thus
$I^\bullet(M)^N$ computes the required $N$-cochains.

Each $I^j(M)^N$ is coinduced as a $Q$-module.  Explicitly,
the function
\[
 F(\overline g_0;r_1,\ldots,r_j)
    =g_0^{-1}f(g_0,g_0r_1,\ldots,g_0r_j)
\]
identifies it with
$C_{\mathrm{cts}}(Q,C_{\mathrm{cts}}(G^j,M))$.
The expression is independent of the lift $g_0$ by
$N$-invariance; $Q$ acts by left translation on its first
variable.  Continuity in both directions follows from
the chosen section.  The positive $Q$-cohomology of this
module vanishes: in the homogeneous bar complex,
evaluation after inserting the identity in the free
$Q$-coordinate is a continuous contracting homotopy.
The degree-zero invariants are $I^j(M)^G$.

Consider the first-quadrant double complex
\[
 C^s_{\mathrm{cts}}(Q,I^t(M)^N).
\]
Taking horizontal cohomology first therefore identifies
its total cohomology with that of $I^\bullet(M)^G$,
namely continuous $G$-cohomology.  Taking vertical
cohomology first uses the stipulated family property,
with parameter $Q^s$, and gives
$C^s_{\mathrm{cts}}(Q,V_t)$.  This proves
\eqref{eq:continuous-HS}.  In each total degree there are
only finitely many bidegrees, so the standard filtration
converges.

All comparisons, prism sums, and contractions use
continuous evaluation, group operations, and finite sums
at one coefficient stage.  For a stage union, perform
them before taking the filtered colimit; this preserves
their identities and exactness.  Replace every parameter
$Q^s$ by $T\times Q^s$ to obtain the parameterized
assertion.  A map that is a quasi-isomorphism for all
these $N$-parameters is an isomorphism on the vertical
cohomology pages, hence on total cohomology.  This proves
the final assertion.
\end{proof}

\subsection{Removing the roots and computing the norm quotient}

Put
\[
 H'=\Nrd^{-1}(\mu_2)\subset P,\qquad
 \chi:H'\twoheadrightarrow\mu_2\hookrightarrow k^\times .
\]
The scalar $-1\in P$ has norm $-1$, and is central, so
$H'=H\times\mu_2$.  Also $P/H'=1+3\mathbb Z_3\cong\mathbb Z_3$.
The norm quotients used below have continuous sections.
For an explicit section on the principal units, choose an
unramified cubic subfield $E\subset D_3$ and
$u\in\mathcal O_E$ with $\operatorname{Tr}_{E/\Qp}(u)=1$.
Such a $u$ exists because the trace of the unramified
residue-field extension is surjective.  Then
\[
 z\longmapsto \exp(u\log z),\qquad z\in1+3\mathbb Z_3,
\]
is a continuous homomorphism into $\mathcal O_E^\times$
whose reduced norm is $z$.  Together with the central
order-two scalar it supplies the sections; thus
Lemma~\ref{lem:continuous-extension} applies.

\begin{theorem}[The ordinary coefficient calculation]
\label{thm:coefficient-calculation}
The natural inclusion $L\hookrightarrow R$ induces, for
every compact profinite $T$, a quasi-isomorphism
\begin{equation}
 \label{eq:ordinary-perfect-P-comparison}
 \Ccts(P,L;T)\xrightarrow{\ \sim\ }\Ccts(P,R;T).
\end{equation}
It is $J$-equivariant and natural in $T$.  With discrete
cohomology groups, evaluation gives
\begin{equation}
 \label{eq:ordinary-four-rows}
 H^a\Ccts(P,L;T)=
 \begin{cases}
 C_{\mathrm{cts}}(T,k),&a=0,1,\\
 C_{\mathrm{cts}}(T,\ell),&a=3,4,\\
 0,&\text{otherwise}.
 \end{cases}
\end{equation}
Here the identifications in degrees one and four use the same nonzero
element of $\Hom_{\mathrm{cts}}(1+3\mathbb Z_3,\Fp)$.
The $J$-action is trivial in degrees zero and one.
On each upper line, its restriction to
$\mathcal O_F^\times$ is the nontrivial norm character
$\delta$.
\end{theorem}

\begin{proof}
The transformation law of the distribution parameter
gives the $H'$-equivariant coefficient sequence
\begin{equation}
 \label{eq:distribution-sign-sequence}
 0\longrightarrow L\longrightarrow R
 \xrightarrow{D}R(\chi)\longrightarrow0.
\end{equation}
Here $R(\chi)$ means that $h\in H'$ acts as
$\chi(h)$ times its action on $R$.  The kernel,
surjectivity, and continuous lifts at finite larger
stages, including arbitrary profinite parameters, were
proved in Lemma~\ref{lem:distribution-operator}.
Thus \eqref{eq:distribution-sign-sequence} is exact on
our $H'$-cochain complexes.

By Theorem~\ref{thm:decompletion}, the two nonzero
$H$-cohomology groups of $R$ have trivial $\mu_2$-action.
After twisting by $\chi$ they have the nontrivial
order-two action.  Averaging over $\mu_2$ is exact
because $2$ is invertible in $k$.  Applying
Lemma~\ref{lem:continuous-extension} to
$H'=H\times\mu_2$, with all profinite parameters, shows
that $\Ccts(H',R(\chi);T)$ is acyclic.
The long exact sequence of
\eqref{eq:distribution-sign-sequence} therefore makes
\[
 \Ccts(H',L;T)\longrightarrow\Ccts(H',R;T)
\]
a quasi-isomorphism for every $T$.
Its map is the natural inclusion, so it retains the
commuting $P$- and $J$-actions although the auxiliary
operator $D$ need not do so.
Apply the comparison assertion of
Lemma~\ref{lem:continuous-extension} to
$H'\subset P$ to obtain
\eqref{eq:ordinary-perfect-P-comparison}.

It remains to calculate the cohomology of $R$.
Exact averaging gives
$H^t\Ccts(H',R)=k$ for $t=0$, $\ell$ for $t=3$,
and zero otherwise.  The quotient
$Q=P/H'\cong\mathbb Z_3$ acts trivially on both
lines by Theorem~\ref{thm:decompletion}.
The two-term completed-group-ring resolution
\[
 0\longrightarrow\Fp[[Q]]
 \xrightarrow{\gamma-1}\Fp[[Q]]
 \longrightarrow\Fp\longrightarrow0
\]
for a topological generator $\gamma$ computes
$H^s(Q,V)=V$ in degrees $s=0,1$ and zero otherwise
when $Q$ acts trivially on $V$.
The spectral sequence \eqref{eq:continuous-HS} thus
has exactly the four terms in
\eqref{eq:ordinary-four-rows}; no differential is
possible, and each total degree contains one term.
The same resolution and the parameterized version of
Lemma~\ref{lem:continuous-extension} give the formula
for arbitrary $T$.

Finally $J$ commutes with $P$ and hence acts trivially on
the quotient $Q$ and its degree-one cohomology line.
The low-degree coefficient line consists of constants.
The upper coefficient line has the action specified
in Theorem~\ref{thm:decompletion}.  This proves all
assertions about the $J$-action.
\end{proof}
